\section{Introduction}\label{sec:intro}

Self-adaptive systems choose actions under uncertainty, and when the stakes are high their controllers must come
with guarantees: a robot must reach its waypoints in order with high probability while rarely meeting a moving
obstacle, and an autonomous underwater vehicle must finish its inspection within an energy budget. Probabilistic
model checking gives such guarantees for one Markov decision process (MDP) at a time. Deployed systems, however,
meet a \emph{family} of MDPs: the map changes, the grid grows, the mission gets longer. Synthesising and checking a
new controller for every variant is expensive at run time, and the controllers it produces (per-state strategy
tables) are hard to inspect. What the engineer wants is one readable controller for the whole family, with a
guarantee on every variant the system will meet.

Large language models (LLMs) are good at writing short, general programs, and recent work uses them to write
generalized plans and policies for classical planning domains~\cite{silver2024generalized,stein2026strategy}.
Their output is tested, not certified, and only on deterministic domains, where running a plan is enough to check
it. Methods for families of MDPs, on the other hand, either need every member to share one state
space~\cite{andriushchenko2024trees,heck2026smpmc} or learn a decision tree over raw state variables from the
optimal strategies of small instances~\cite{azeem2025onetwothreego}, which cannot follow quantities that move with
the instance (a goal's position, the length of a game).

We combine the two. An LLM writes an ordered list of rules (\emph{first match wins}) over a vocabulary the domain
provides: the state variables, named instance constants (the grid size, the goals' coordinates) and features such as
``this move hits a wall''. A loop checks the rules with PRISM~\cite{kwiatkowska2011prism} on a few small
\emph{training} instances at once and feeds back which requirements fail, where probability is lost, and which rules
are responsible. The rule list then is frozen and certified on held-out instances, larger ones included: a rule list
that leaves some states open, or allows several actions in a state, is checked against \emph{every} way of
resolving those choices, so a passing check holds for any deployment of the rules.

The LLM writes the strategy, but the verifier does more than check it. Rules that allow several actions (``move
towards the goal unless it is a wall'') are convenient to write, yet the worst case lets an adversary pick among the
allowed moves, and rules tuned against that adversary do not carry over. We therefore let the verifier
\emph{refine} the rules: in every round it resolves each such choice into the order of actions that does best on the
training instances, and the LLM sees how the resulting deterministic policy fails. This step costs model checking
only, needs no optimal strategies, and turned out to be what makes the rules transfer.

\paragraph*{Contributions}
\begin{itemize}
  \item A loop in which an LLM writes one rule list for several training instances of an MDP family, refined by
        verification against several probabilistic requirements and certified on every held-out instance
        (Section~\ref{sec:approach}).
  \item Verifier-side refinement: choices left open by the LLM are resolved, every round, into the action order that
        verification finds best on the training instances (Section~\ref{sec:approach-refine}).
  \item An evaluation on gridworld (random grids and three parameterized families) and on the Pac-Man benchmark of
        the Quantitative Verification Benchmark Set~\cite{hartmanns2019qcomp}, against decision trees learned from
        optimal strategies (as in 1--2--3--Go!, with and without our features), PPO~\cite{schulman2017ppo}, a search
        over the same rule language, and a hand-written controller (Sections~\ref{sec:evaluation}--\ref{sec:results}).
\end{itemize}

\todo{one paragraph with the headline numbers once the final runs are in}
